% BANKS-SOURCE: pdf=249 print=239 kind=problem-section
\begin{banksproblems}

% BANKS-SOURCE: pdf=249 print=239 kind=problem id=10.1
\BanksProblem{10.1}{*}
Discuss vacuum decay for a scalar field theory in $D$ dimensions. That is, we
have a single scalar field, with the same kind of 2-minimum potential we
discussed in the case of quantum mechanics. The Euclidean path integral is over
all field configurations $\phi(\mathbf x)$ that approach $\phi_f$, the higher minimum
of the potential, at infinity (set the potential equal to zero there). Assume
that the instanton configuration is $SO(D)$-invariant. Find the field equation
for it. If you call the Euclidean radius $t$, you will find that the equation is
identical to a mechanics problem in the upside-down potential, with a
time-dependent friction term $(D-1)/t$.

% BANKS-SOURCE: pdf=250 print=240 kind=prose
Using simple energetic considerations, argue that there is a finite-action
bounce solution. Write down the equations for small fluctuations around the
bounce. Since the bounce is spherically symmetric, they can be separated by
expanding in $D$-dimensional spherical harmonics. How many exact zero modes of
the solution exist, following from symmetries? In what representation of $SO(D)$
do they transform? Try to argue that there is exactly one negative-eigenvalue
fluctuation mode. (Don't try too hard, the general result for a large class of
quantum systems can be found in \cite{banks-ref-151}. If you can do it easily
for this special case, do it.) The answer to most of this exercise can be found
in the Coleman book \cite{banks-ref-138}. Try to do it yourself, but you can go
there for hints if you get stuck.

% BANKS-SOURCE: pdf=250 print=240 kind=problem id=10.2
\BanksProblem{10.2}{*}
Prove Derrick's theorem. Consider a Lorentzian scalar theory with mostly-plus
signature whose static energy has gradient and potential terms
$(\nabla\phi^i)^2+V(\phi)$ in $d_{\rm s}>2$ spatial dimensions. Consider any finite-energy non-singular field
configuration and show that you can lower its energy by considering
$\phi(ax)$ with appropriate $a$. This shows that, without further constraints,
there are no stable solitons in such a theory. Generalize this to an arbitrary
number of scalars. Now consider a complex scalar field with a charge quantum
number. The potential is a function of $\Phi^*\Phi$. For a
$D$-dimensional Lorentzian spacetime, show that the positive-frequency fields
of the form $\Phi(t,\mathbf x)=\ee^{-\ii\omega t}\phi(\mathbf x)$, with
$\omega>0$, carry an electric charge

% BANKS-SOURCE: pdf=250 print=240 kind=display
\[
  Q=\ii\int\dd^{D-1}\mathbf x
  \left(\Phi^*\partial_t\Phi-\Phi\partial_t\Phi^*\right).
\]

Evaluate the energy of such configurations. Show that Derrick's theorem no
longer applies. In the appendix to the chapter on lumps in Coleman \cite{banks-ref-138},
you will find an example of a theory of this type that has stable charged
solitons. Such objects are called Q-balls.

% BANKS-SOURCE: pdf=250 print=240 kind=problem id=10.3
\BanksProblem{10.3}{*}
Consider the quantum Lagrangian

% BANKS-SOURCE: pdf=250 print=240 kind=display
\[
  L=\frac12(\dot\phi)^2-V(\phi),
\]

where $V$ is periodic with period $2\pi$. There are two different
interpretations of this system. In the first, $\phi$ lives on a circle. In the
second, $\phi$ lives on the real line, in the presence of a periodic potential.
In the second interpretation the system is invariant under translation of $\phi$
by $2\pi$ and contains a global symmetry operator $T$ that commutes with the
Hamiltonian. We can diagonalize the Hamiltonian in the sector satisfying

% BANKS-SOURCE: pdf=250 print=240 kind=display
\[
  T\psi(\phi)=\psi(\phi+2\pi)=\ee^{\ii\theta}\psi(\phi),
\]

where $\theta$ is called the Bloch momentum. In the first interpretation,
translating $\phi$ by $2\pi$ does nothing, and wave functions are required to be
periodic. In either case, we are allowed to add the term
$\delta L=\theta\dot\phi$ to the Lagrangian, which is a total derivative,
violating time-reversal invariance, which is analogous to the $F\widetilde F$
term in gauge theories. Show that, in either interpretation, there are
finite-action Euclidean instantons that have value $2\pi$ times an integer for
this topological term. Show that, in the case that $\phi$ lives on the real
line, the path integral with non-zero $\theta$ corresponds to computing the
ground-state expectation values in the lowest energy state with Bloch momentum
$\theta$.

% BANKS-SOURCE: pdf=251 print=241 kind=problem id=10.4
\BanksProblem{10.4}{*}
The instanton solutions of a quantum-mechanics problem are solitonic particle
solutions of field theory with the same Lagrangian in $D=2$ $(1+1)$ dimensions.
Show that, in any higher dimension, the same solutions can be interpreted as
\emph{domain walls} separating two ground states of the theory. Similarly, show
that the instantons of the $D=2$ $(1+1)$-dimensional Higgs model can be thought
of as particles in the $D=3$ $(2+1)$-dimensional version of the model and as
strings or (Abrikosov--Nielsen--Olesen) vortices in the $D=4$ $(3+1)$-dimensional
version. This
problem requires no additional calculation.

% BANKS-SOURCE: pdf=251 print=241 kind=problem id=10.5
\BanksProblem{10.5}{*}
Prove that the instanton solution of the $D=4$ Euclidean gauge theory is invariant
under rotations, in the sense that all gauge-invariant functions of the field are
rotation-invariant.

\end{banksproblems}
